\documentclass[12pt]{article}

% chktex-file 1
% chktex-file 3

\usepackage[utf8]{inputenc}
\usepackage[T1]{fontenc}
\usepackage{amsmath, amssymb}
\usepackage{graphicx}
\usepackage{tikz}
\usepackage{geometry}
\usepackage{booktabs}

\geometry{margin=2.5cm}
\begin{document}

\section*{Actividad de Taller de Simulación de Sistemas}
\addcontentsline{toc}{section}{Actividad de Taller de Simulación de Sistemas}

Deberá estructurar la entrega de su trabajo de acuerdo a los siguientes puntos:
\begin{enumerate}
    \item Modelo de simulación teórico de los problemas asignados.
    \item La simulación implementada en Excel: corresponde a la simulación manual y los correspondientes gráficos.
    \item El programa Java de simulación: El punto de partida es el modelo teórico de Simulación.
    \item Llevar a cabo la validación correspondiente en el modelo de simulación manual y aquel implementado.
    \item Presentar los requerimientos respecto de informe, diapositivas y código fuente ya explicados en clase.
\end{enumerate}

A continuación los ejercicios que hacen a su actividad de taller de simulación.

\section{Parte 1: Dos problemas distintos}

\subsection{Problema 1: Función dada (aplicar transformada inversa)}

La función de densidad proporcionada es:

\[
f(x)=
\begin{cases}
\dfrac{(x-3)^2}{18}{}, & 0 \leq x \leq 6 \\
0, & 0 > x > 6
\end{cases}
\]

El objetivo consiste en generar valores aleatorios mediante el método de transformada inversa y verificar que los valores obtenidos pertenezcan al intervalo $[0,6]$.

\subsection{Problema 2: Gráfico dado (repetición según apuntes)}

En este segundo problema de la Parte 1 se trabaja con el gráfico proporcionado por el docente, siguiendo el mismo procedimiento explicado en los apuntes de clase.

\begin{center}
\begin{tikzpicture}[scale=1.2]

% Ejes
\draw[->] (0,0) -- (8,0) node[right] {$x$};
\draw[->] (0,0) -- (0,4) node[above] {$f(x)$};

% Triángulo
\draw (1.5,0) -- (4,3) -- (6.5,0);

% Línea punteada
\draw[dashed] (4,0) -- (4,3);

% Etiquetas
\node[below] at (1.5,0) {$a$};
\node[below] at (4,0) {$b$};
\node[below] at (6.5,0) {$c$};

\end{tikzpicture}
\end{center}

\section{Parte 2: Aplicaciones de Simulación}

\subsection{Problema 1}

La demanda diaria y el tiempo de entrega de un cierto producto, siguen las siguientes distribuciones de probabilidad:

\begin{center}
\begin{tabular}{cc}
\toprule
Demanda diaria & Probabilidad \\
\midrule
0 & 0.04 \\
1 & 0.06 \\
2 & 0.10 \\
3 & 0.20 \\
4 & 0.30 \\
5 & 0.18 \\
6 & 0.08 \\
7 & 0.03 \\
8 & 0.01 \\
\bottomrule
\end{tabular}
\qquad
\begin{tabular}{cc}
\toprule
Tiempo de entrega (días) & Probabilidad \\
\midrule
1 & 0.25 \\
2 & 0.50 \\
3 & 0.20 \\
4 & 0.05 \\
\bottomrule
\end{tabular}
\end{center}

La información con respecto a los costos relevantes es la siguiente:

\begin{itemize}
    \item Costo de ordenar: \$50/orden
    \item Costo de inventario: \$26/unidad/año
    \item Costo de faltante: \$25/unidad
\end{itemize}

Si el inventario inicial es de 15 unidades, ¿determine la cantidad óptima a ordenar (q) y el nivel óptimo de reorden (R)? (Asuma que se trabajan 260 días en el año.)

\subsection{Problema 2}

La demanda diaria y el tiempo de entrega de un cierto producto, siguen las siguientes distribuciones de probabilidad:

\begin{center}
\begin{tabular}{cc}
\toprule
Demanda diaria & Probabilidad \\
\midrule
25 & 0.02 \\
26 & 0.04 \\
27 & 0.06 \\
28 & 0.12 \\
29 & 0.20 \\
30 & 0.24 \\
31 & 0.15 \\
32 & 0.10 \\
33 & 0.05 \\
34 & 0.02 \\
\bottomrule
\end{tabular}
\qquad
\begin{tabular}{cc}
\toprule
Tiempo de entrega (días) & Probabilidad \\
\midrule
1 & 0.20 \\
2 & 0.30 \\
3 & 0.25 \\
4 & 0.25 \\
\bottomrule
\end{tabular}
\end{center}

Si el producto no está disponible cuando es requerido, el cliente puede esperar la llegada de un nuevo lote por un tiempo limitado, es decir, si el cliente decide esperar 2 días y la mercancía no llega en ese tiempo, entonces, la demanda de este cliente se considera perdida. La distribución de probabilidad del tiempo que un cliente está dispuesto a esperar para que se le surta su pedido, es la siguiente:

\begin{center}
\begin{tabular}{cc}
\toprule
Tiempo de espera (días) & Probabilidad \\
\midrule
0 & 0.40 \\
1 & 0.20 \\
2 & 0.15 \\
3 & 0.15 \\
4 & 0.10 \\
\bottomrule
\end{tabular}
\end{center}

La información con respecto a los costos relevantes es la siguiente:

\begin{itemize}
    \item Costo de ordenar: \$100/orden
    \item Costo de inventario: \$52/unidad/año
    \item Costo de faltante suponiendo que el cliente espera: \$20/unidad
    \item Costo de faltante suponiendo que el cliente NO espera: \$50/unidad
\end{itemize}

Si el inventario inicial es de 100 unidades, ¿determine la cantidad óptima a ordenar (q) y el nivel óptimo de reorden (R)? (Asuma que se trabajan 260 días en el año.)

\end{document}